\documentclass[11pt,a4paper]{ltjsarticle}
\usepackage[margin=25mm]{geometry}
\usepackage{amsmath,amssymb,booktabs,longtable,array}
\usepackage[haranoaji]{luatexja-preset}
\usepackage[unicode,hidelinks]{hyperref}
\usepackage{xurl}
\newcommand{\Gr}{\operatorname{Gr}}
\newcommand{\PGL}{\operatorname{PGL}}
\title{nc-p3研究方針と再開手順}
\author{}
\date{2026年11月15日\quad 第16版}
\begin{document}
\maketitle

\section{最初に読む：目標と達成基準}
標数零の代数閉体上、標準14641型の四次元quadratic Koszul AS正則
$\mathbb Z$-代数を\textbf{添字保持同型}で完全分類する。
全対象にdomain、noetherian、PBW、CYを追加仮定しない。
$qgr$同値、導来同値、添字移動の商、周期構造を選んだ同型は別問題とする。
達成基準は、正則locusの必要十分条件、明示的な族による全被覆、特殊値を含む同型判定、
全周期層、および既約成分・次元・特殊化の決定である。完全分類は未完成である。

直近の数学的目標は、graded同型分類とindexed分類の次元差を、
twist忘却による減少と非1周期的な対象の出現に分けて決めること。
前者は一般量子成分で決定した。後者について一般最小周期2の九次元の同型類像を構成したが、
その像が成分を尽くすか、Clifford点の成分次元がいくつかは未決定である。

正本はGitHub \texttt{uedakazushi/nc-p3}の\texttt{main}。
再開時は先端を確認し、本方針、\path{AGENTS.md}、主論文\path{nc-p3.tex}を読み、
必要な計算・障害は\path{research/notes_ja.tex}へ進む。
証明位置と課題は\path{research/STATE.json}、文献と外部依存は\path{research/BASELINE.md}、
今回の生成・計算・組版・保存の別状態は\path{research/validation.json}に置く。
GitHubへのpush成功時はLibraryへの重複保存を行わない。

\section{現状評価と主要な到達点}
公表原稿の中心は\textbf{非1周期的な代数的族とその幾何}である。
有限四窓表示は分類の基礎、量子成分とClifford点の比較は変形の説明、
境界退化は像の大きさと接空間の違いを示す具体例として位置づける。
以下の「証明」は原稿に論証があるという意味であり、形式化・第三者査読・新規性確定とは区別する。

\begin{longtable}{p{.24\linewidth}p{.69\linewidth}}
\toprule
主題 & 確定範囲と証明位置\\
\midrule
\endhead
標準4周期と有限表示 &
Frobenius traceから標準4周期を構成し、幾何学的同型類を四窓へ還元。
AS正則性が既に認証された標準四窓の同型・周期は212未知数388式で判定する。
双対Frobenius性の有限検査はKoszul判定を含まない。
\path{sections/four_window.tex}。\\
量子成分 &
一般点でgraded同型類6次元、indexed同型類3次元。
三つのtwist方向を忘れる比較である。
\path{sections/local_moduli.tex}の
\texttt{thm:quantum-orbits}, \texttt{thm:quantum-local}。\\
Clifford点 &
選んだgraded表示の一次変形は9、indexed一次変形は15。
全15変数の二次障害は下記$Q$。
\texttt{thm:clifford-tangent}, \texttt{thm:clifford-quadratic}。
成分次元14は予想にとどまる。\\
七・九変数族 &
七変数族は全$\mathbb A^7$上PBW自由。
九変数族の基底は$S_9=\operatorname{Spec}k[u,v,r,s,a,b,c,t,h,(1+th)^{-1}]$、
$t=h=0$が七変数族。ともに全ファイバーAS正則、同型類像の次元は7と9、
一般最小周期2。
\path{sections/algebraic_family.tex},
\path{sections/elliptic_ore_family.tex}。\\
Ore表示と不変量 &
固定三生成Clifford部分上の次数一導分は通常6、符号捩れ9。
基底変更を除く九変数正規形は指定Ore構成の一般表示を尽くす。
楕円曲線・2等分点は内在的で、$U(th)$は同型の必要不変量。
\path{sections/ore_complete.tex},
\path{sections/isogeny_invariant.tex}。\\
境界チャート &
五変数表示の同型類像は2次元、一般最小周期2。
別の三生成部分を用いるとOre表示を持ち、七変数族の軌道閉包に含まれる。
\path{sections/nonore_family.tex},
\path{sections/nonore_period_two.tex},
\path{sections/ore_boundary.tex}。\\
閉包の特異性 &
位置0の関係空間の軌道閉包$Z_7,Z_9\subset\Gr(6,16)$について
$\dim Z_7=37,\ \dim Z_9=39$、
Clifford点の接次元はそれぞれ39以上、41以上。
\path{sections/ore_boundary.tex},
\path{sections/nine_closure.tex}。\\
\bottomrule
\end{longtable}

再開時に使う二つの式は
\[
 Q=o_0o_5+o_1o_2-o_2o_3+o_3o_4,\qquad
 U(x)=\frac{4096(x+1)^3}{x(x-8)^3},\quad x=th.
\]
$Q$は全変形環ではなく二次障害であり、$U$は十分な同型不変量ではない。
境界を加えた11個の個別に積分可能な接方向では$Q|_L=\lambda c$となるため、
その全方向を接像に持つ滑らかな基底上の同時積分は不可能である。
閉包の特異性は、全モジュライの特異性や新しい成分を意味しない。

\section{未解決の核心と優先順位}
\paragraph{最優先：ALGEBRAIC-FAMILY-008。}
九変数族の\emph{一般点}の変形を調べ、像の最大性を判定する。
Clifford特殊点だけの接空間計算を重ねるより、既存族が一般点で尽くす方向と不足方向を分ける方が
成分問題に直結する。明示的な代数的変形を優先する方針は維持する。
着手は次の順に行う。
\begin{enumerate}
\item $th(1+th)(th-8)\ne0$かつ12点が単純な開集合で、
周期2および標準周期4の変形複体を構成する。
階数を標本で探索した後、記号的小行列式または構造論で一般値を証明する。
現在の$\mathrm{PBW}$書換えコードを一般$t,h$へそのまま流用せず、
中心平方上のClifford基底とOre積を使う。
\item 九パラメータと基底変更の接像を全変形と比較する。
全射なら滑らかな源を用いた局所完全性の議論を行う。
余核があればその点の障害写像（Clifford点では$Q$）と必要な高次障害で候補を絞り、
実際の代数的族を構成する。
接次元が大きいことだけから高次元成分を宣言しない。
\item 偏極付き楕円曲線、2等分点、12点因子と第四座標の持上げデータから、
indexed同型の必要十分条件を求める。
$U(th)$で高々4候補へ減らした後も残る有限同一視を調べる。
12点の支持だけで関係空間が回収されるとは未証明なので、必要なら非被約点の局所構造も含める。
\end{enumerate}
五変数境界族の七変数族への閉包包含と一般最小周期2は解決済みであり、再探索しない。
未解決なのは全AS正則locusにおける成分性・最大性と、閉包の外へ出る代数的変形である。

\paragraph{次順位：PERIOD-ROOTS-007。}
Clifford点の射影自己同型と全1周期構造を分類する。
特定のgraded表示の9次元接像だけから1周期的部分全体を決定したことにはならない。
標準四窓の有限同型式を使い、自己同型と周期構造の根を分離する。
この比較が終わるまで非1周期的部分の余次元を確定しない。

\paragraph{補助手段：MODULI-HIGHER-006。}
二次障害$Q$を候補選別に用い、代数的族の構成または不存在証明に必要な高次補正だけを計算する。
一つの補正の失敗は非持上げの証明にならず、有限個の三次サンプルの成功も全次数を保証しない。

\paragraph{全体目標への接続：Z4-AS-004。}
標準四窓への抽象的な全対象還元は得られている。
残るのは一般四窓のKoszul/AS条件の明示的な必要十分判定、
正則な代数的チャートによる全被覆、非PBW点・閉特殊値・最小周期4の層と族間同型である。
相対族の大域的下降、全成分の次元と特殊化も未解決として保持する。
元の完全分類目標を局所例の構成で達成済みにしない。

\section{公表原稿と文献比較}
主論文は当面一つに保ち、有限表示、局所比較、代数的族、同型不変量、境界の順に読む構成とする。
七変数族の独立した全パラメータ証明や境界の有用な証明を削らない。
研究ノートには本文にない基底データ、反例、失敗した方法の障害を残し、本文の証明を重複させない。

新規性比較の優先対象は、Vancliff--Van Rompay--Willaertのgraded Clifford/Ore変形、
Stephenson--Vancliffの三生成Clifford部分を用いた構成、
Bhatoy--Ingalls--LaRoche--Nookalaのgraded変形計算、
Bhatoy--Ingallsのsuperpotential不変量である。
本稿の候補となる寄与は添字依存の交代変形と非1周期性、同型類像の次元、
境界の具体的な幾何に置く。
一般的なClifford/Ore手法そのものや、graded側の既知の9次元を新規成果として売らない。
確認した一次資料・参照位置と未確認点は\path{research/BASELINE.md}に集約する。

\section{再現・検証と継続}
全コマンドの正本は\path{README.md}の「再現」節。
Python/SymPyによる構造検査と厳密検算、LuaLaTeXによる主論文・方針・ノートの組版を行う。
証明書は現行コードと入力から再生成し、数学的内容と環境表示の差を分けて比較する。
\path{scripts/check_project.py}の成功は数学的証明を意味しない。
未実行の検証を過去版のPASSで置き換えない。

研究の区切りでは本文・要約・方針・索引の不整合を直し、再開可能な現行一式をmainへ保存する。
不要な作業履歴・二重の先端表・旧版の未解決項目は蓄積しない。
保存結果は遠隔を読み戻して確認する。外部as4の結果は固定した出典を使い、
その全計算を再実行していない場合は明記する。

\section{常設の研究原則と保存許可（具体的計画とは独立）}
数学的重要性、新規性、理解の深化、応用、正確性、実現可能性を総合し、目標と方針を柔軟に最適化する。
重要な変更では有効な理由、元の目標の達成・未達を簡潔に整理し、履歴を積み重ねない。
証明済み、追加仮定付き、未完、予想、計算的観察を区別する。
原稿に証明があること、今回監査したこと、計算で確認したこと、形式化・査読は別状態とする。
誤りがあれば依存する本文・要約・方針を直し、有用な反例と障害だけをノートに保持する。

公表論文は原則日本語LuaLaTeXで書き、関連する一次文献の仮定・適用範囲と寄与をIntroductionで説明する。
未確認の書誌・内容を推測しない。公表価値のある既存の主張と十分な証明を現行論文に保持する。
ノートは作業日誌ではなく、未収載の有効な知識、反例、計算、失敗した方法の障害、文献比較を整理する。
論文の完全な証明はノートへ重複させず参照する。

再開に必要な内容を現行の方針・ノート・論文・最小限のコードと入力に集約する。
旧bundle、旧原稿、過去会話を再開の前提にしない。
旧ZIP、バックアップ、変更履歴、作業ログ、キャッシュ、コンパイル中間物、重複資料をbundleから除く。
数学的内容・証明書・再取得困難な入力を容量制限のために削らない。
日付・版・参照・記法を整合させ、更新は現行内容の整理・置換によって行う。

ユーザーは研究の区切りごとの承認不要の保存を許可している。
本プロジェクトの正本はGitHub \texttt{uedakazushi/nc-p3}の\texttt{main}とし、通常のfast-forwardで直接保存する。
PR、force push、GitHub Actionsは使わず、別リポジトリas4へ無断で書き込まない。
GitではPDFを追跡せず、ソースと再現可能な資産を追跡する。
GitHub保存が成功した場合はLibraryへ重複保存せず、GitHubへ保存できない場合は常設許可によりLibraryへの緊急退避を試みる。
既存Libraryファイルは明示指示なく削除しない。
操作手段がなければ保存済みと報告せず、生成・検証・GitHub保存・Library保存を別々に明示する。
配布前に欠落・参照切れ・組版・主要計算・ZIP整合性を確認し、未実行の検証や再現上の制約を明記する。
\end{document}
